\section{What We Propose Next}
\label{sec:next}

\subsection{Two requests, and what each unblocks}

\paragraph{A seeded-defect build.} A version of any application you choose,
with a known list of injected faults, and nothing else. With it we can report
precision and recall for the first time, and so can the field: no reviewed
system reports either. Without it, every defect count this system produces
remains a candidate list whose accuracy is unmeasured, and no return-on-investment
claim about the tool can be made honestly, including ours.

This is the single highest-value thing we could be given. It requires no
integration work on your side beyond producing the build and the ground-truth
list.

\paragraph{A defect export.} An extract from a real issue tracker, with
severity, area and root cause, for an application we can also run against.
ETA-REQ-301 is built and works on synthetic data; it has never influenced test
selection using real history. With an export we can report whether
defect-weighted retrieval actually directs the agent toward the areas that
historically break, which is the claim the requirement was written to support.

\subsection{Engineering we would do next}

In the order we would take them, each justified by something in
\Cref{sec:limits} rather than by novelty.

\begin{enumerate}[leftmargin=2.0em,itemsep=4pt]
  \item \textbf{Content-independent state signatures.} Replace whole-tree
    comparison with a signature over structure alone, so a revisited screen is
    recognised after its contents change. This directly attacks the autonomy
    ceiling on targets without stable identifiers.
  \item \textbf{Extraction metering.} An output ceiling, a reduced sample count
    once extraction stabilises, and a cache that refuses to re-read an
    unchanged document. Ingestion is the largest cost line and the easiest to
    reduce.
  \item \textbf{Parallel device execution.} The architecture already separates
    planning from execution across processes, so one planner can drive several
    devices. This turns a seven-hour campaign into
    roughly a one-hour one, and wall clock is the binding constraint on how
    often a campaign can run.
  \item \textbf{Provisional specification generation.} A first exploratory pass
    that produces requirements the later passes can target, for the case where
    no specification exists. This is what would make DR-01 true in intent and
    not only in letter.
  \item \textbf{Deduplication against elements rather than text.} Removes the
    over-counting described in \Cref{sec:limits}.
  \item \textbf{A controlled comparison of the two planners.} The tool-calling
    planner is better formed and is switched off because no evidence says it
    produces better campaigns. One controlled pair of campaigns settles it.
\end{enumerate}

\subsection{Where this could go inside an organisation}

The system is operable by one engineer from a dashboard, which was a
requirement rather than an accident. Its natural place is a nightly run against
a build, with a QA engineer triaging candidates in the morning. The cost of the
inference is negligible at that cadence; the cost that matters is the
engineer's time spent confirming candidates, which is why precision, once it
can be measured, is the number that decides whether the tool is worth
operating.

Nothing in the architecture assumes a single application or a single platform.
Adding a target is a configuration file, and the evidence for that claim is
five targets across two platforms with no code change between them.

\vspace{10pt}
\noindent\rule{\textwidth}{0.4pt}

\vspace{4pt}
\noindent\footnotesize The repository, the campaign records and the 123
trajectory directories underlying every figure in this report are available to
Samsung Research and Development Institute Bangladesh on request.
